\documentclass[conference]{IEEEtran}

% ------------------------------------------------
% Packages
% ------------------------------------------------
\usepackage{amsmath, amssymb}
\usepackage{graphicx}
\usepackage{cite}
\usepackage{url}
\usepackage{booktabs}
\usepackage{float}
\renewcommand{\refname}{Bibliografía}
% ------------------------------------------------
% Document
% ------------------------------------------------
\begin{document}

% ------------------------------------------------
% Title
% ------------------------------------------------
\title{Evaluación de la Vulnerabilidad de la Red Vial Primaria de Carga de Colombia Mediante Teoría de Grafos}

% ------------------------------------------------
% Authors
% ------------------------------------------------
\author{
    \IEEEauthorblockN{
        David Catral,
        Gabriel Salcedo,
        Isaac Suarez
    }
    \IEEEauthorblockA{
        Escuela de Economía y Negocios\\
        Universidad del Rosario\\
        Bogotá, Colombia\\
\{jose.catral, gabriel.salcedo, isaac.suarez\}@urosario.edu.co    }
}

\maketitle


% ------------------------------------------------
% I. Introducción
% ------------------------------------------------
\section{Introducción}
Una red de transporte se representa de forma natural como un
grafo ponderado: los vértices son localidades o cruces, las aristas son los
tramos viales que los conectan y el peso de cada arista es el costo de
recorrerla. Sobre esa representación, la teoría de grafos ofrece herramientas
directas para estudiarla. El algoritmo de Dijkstra \cite{dijkstra1959} calcula
la ruta de mínimo costo entre dos puntos; la centralidad de intermediación
\cite{brandes2001} mide qué componentes concentran el tránsito de la red; la
eficiencia global \cite{latora2001} cuantifica cuánto se degrada el sistema al
perder uno de ellos; y la búsqueda en profundidad \cite{tarjan1972} identifica
los puentes, es decir, las aristas cuya remoción desconecta el grafo.
 
En Colombia estas herramientas tienen una aplicación inmediata. La geografía
montañosa y la distancia entre los centros de producción, los puertos y los
principales mercados hacen que el transporte de carga dependa de un número
reducido de corredores troncales y transversales, que conforman la red vial
primaria \cite{invias_opendata}. Cuando uno de esos tramos se cierra por
derrumbes, crecientes, bloqueos u obras, la carga debe tomar una ruta alterna
más larga y más costosa, y en algunos casos no existe ruta alterna
\cite{cambio_radiografia}. Saber de antemano cuáles tramos producen el mayor
daño permitiría priorizar el mantenimiento y la inversión sobre criterios
objetivos.
 
Este proyecto modela la red vial primaria como un grafo ponderado, le asigna los
flujos de carga reportados en el Registro Nacional de Despachos de Carga por
Carretera (RNDC) \cite{rndc}, identifica los tramos críticos y estima el
sobrecosto de las rutas alternas que resultarían de su cierre. El enfoque sigue
la línea de trabajo sobre vulnerabilidad de redes viales, que mide la
importancia de un tramo por el aumento de costo que produce su interrupción
\cite{jenelius2006}, \cite{scott2006}. El resto del documento está organizado de
la siguiente manera. La Sección~\ref{sec:descripcion} describe el problema y la
Sección~\ref{sec:objetivos} presenta el objetivo general y los objetivos
específicos. Este documento corresponde a la entrega inicial del proyecto; el
marco teórico, el modelamiento del problema y la solución propuesta se
desarrollan en las entregas posteriores.
% ------------------------------------------------
% II. Descripción del Problema
% ------------------------------------------------
\section{Descripción del Problema}
\label{sec:descripcion}

El transporte por carretera moviliza la mayor parte de la carga interna del
país. Solo en el año 2015, el RNDC registra cerca de $6$ millones de viajes y
$85{,}7$ millones de toneladas despachadas entre municipios \cite{rndc}. Toda
esa carga circula sobre una red primaria que sufre cierres frecuentes por
eventos naturales, bloqueos y obras de mantenimiento \cite{cambio_radiografia}.
Cada cierre obliga a desviar los vehículos por rutas más largas, lo que aumenta
la distancia recorrida, el tiempo de viaje y el flete pagado, y ese sobrecosto
termina reflejándose en el precio final de los bienes transportados.
 
El problema es que no resulta evidente cuáles tramos son realmente críticos. Un
tramo puede ocupar una posición central en la red y sin embargo transportar poca
carga, porque la demanda de transporte no se reparte de manera uniforme sobre el
territorio. Otro puede mover un volumen alto y tener una vía paralela de costo
similar, de modo que su cierre casi no afecta al sistema. Lo que hace crítico a
un tramo es la combinación de dos cosas: cuánta carga lo utiliza y qué tan buena
es su mejor alternativa. Responder esa pregunta exige un modelo que integre la
estructura física de la red con el uso real que se hace de ella, y la teoría de
grafos permite construirlo de forma cuantitativa y reproducible, en lugar de
depender de percepciones sobre cuáles vías son importantes.
 

% ------------------------------------------------
% III. Objetivos (General y Específicos)
% ------------------------------------------------
\section{Objetivo General y Objetivos Específicos}
\label{sec:objetivos}
\subsection{Objetivo General}
Evaluar la vulnerabilidad de la red vial primaria de carga de Colombia mediante
teoría de grafos, identificando sus tramos críticos y cuantificando el
sobrecosto de transporte de las rutas alternas que resultarían de su
interrupción.

\subsection{Objetivos Específicos}
Pasos técnicos y secuenciales necesarios para alcanzar el objetivo general:
\begin{itemize}
    \item Construir la red vial primaria nacional como un grafo ponderado a
    partir de la cartografía oficial de INVIAS y caracterizar su estructura:
    orden, tamaño, componentes conexas, puentes y vértices de corte.
 
    \item Construir la matriz origen--destino de carga a partir del RNDC y
    asignarla sobre el grafo mediante rutas de costo mínimo, validando las
    distancias estimadas contra los kilómetros reportados en los despachos.
 
    \item Identificar los tramos críticos de la red contrastando su
    centralidad topológica con la carga que efectivamente los atraviesa.
 
    \item Simular el cierre de dichos tramos y cuantificar el sobrecosto de
    transporte de la ruta alterna, así como el volumen de carga que quedaría sin
    ruta disponible.
\end{itemize}

% % ------------------------------------------------
% % IV. Marco Teórico
% % ------------------------------------------------
% \section{Marco Teórico}
% Introduce los conceptos fundamentales de la teoría de grafos requeridos para entender tu proyecto.

% \subsection{Conceptos Fundamentales}
% Define matemáticamente los componentes del modelo. Por ejemplo, un grafo se puede representar formalmente como:
% \begin{equation}
% G = (V,E)
% \end{equation}
% donde $V$ representa el conjunto de vértices y $E$ representa el conjunto de aristas.

% \subsection{Algoritmos o Métodos Utilizados}
% Explica conceptualmente los teoremas o algoritmos específicos que aplicarás más adelante (ej. Algoritmo de Dijkstra, componentes conectadas, centralidad, etc.).

% % ------------------------------------------------
% % V. Modelamiento del Problema
% % ------------------------------------------------
% \section{Modelamiento del Problema}
% Explica detalladamente cómo traduces el problema real de la Sección II a un modelo matemático de grafos basado en la Sección IV.

% Abstrae los elementos del problema real de la siguiente forma:
% \begin{itemize}
%     \item \textbf{Vértices ($V$):} Representan [definición, ej. routers, ciudades, usuarios].
%     \item \textbf{Aristas ($E$):} Representan [definición, ej. cables de fibra, carreteras, relaciones de amistad].
%     \text{\textbf{Pesos de las aristas:}} Representan [definición opcional, ej. latencia, distancia en km].
% \end{itemize}

% % ------------------------------------------------
% % VI. Solución Propuesta
% % ------------------------------------------------
% \section{Solución Propuesta}
% Describe aquí la estrategia teórica, lógica o algorítmica elegida para resolver el problema ya modelado. Explica la arquitectura de tu propuesta, pseudocódigo o las reglas del sistema que vas a diseñar.

% % ------------------------------------------------
% % VII. Implementación de la Solución
% % ------------------------------------------------
% \section{Implementación de la Solución}
% Detalla el desarrollo práctico real de tu propuesta. Explica las herramientas de software empleadas (ej. Python con NetworkX, MATLAB, C++), cómo procesaste los datos, librerías críticas o detalles de hardware si los hubiera.

% % ------------------------------------------------
% % VIII. Resultados Obtenidos
% % ------------------------------------------------
% \section{Resultados Obtenidos}
% Presenta de forma cuantitativa y cualitativa los hallazgos de tus experimentos, simulaciones o ejecuciones de código.

% % Ejemplo de Figura (Las leyendas van ABAJO en formato IEEE)
% \begin{figure}[H]
%     \centering
%     \includegraphics[width=0.8\columnwidth]{graph.png} % Asegúrate de tener esta imagen en Overleaf o comenta esta línea
%     \caption{Representación en grafo del modelo propuesto.}
%     \label{fig:graph}
% \end{figure}

% Como se observa en la Fig.~\ref{fig:graph}, la topología de la red muestra las propiedades clave analizadas.

% % Ejemplo de Tabla (Las leyendas van ARRIBA en formato IEEE con números romanos)
% \begin{table}[H]
%     \centering
%     \caption{Propiedades del Grafo de la Red}
%     \label{tab:properties}
%     \begin{tabular}{lc}
%         \toprule
%         Propiedad & Valor \\
%         \midrule
%         Número de vértices & 10 \\
%         Número de aristas & 15 \\
%         Grado promedio & 3.0 \\
%         Diámetro de la red & 4 \\
%         \bottomrule
%     \end{tabular}
% \end{table}

% En la Tabla~\ref{tab:properties} se resumen las métricas estructurales obtenidas tras la ejecución del algoritmo sobre la solución implementada.

% % ------------------------------------------------
% % IX. Conclusiones
% % ------------------------------------------------
% \section{Conclusiones}
% Resume los logros principales del proyecto basándote en los resultados. Explica directamente cómo el uso de la teoría de grafos permitió solucionar el problema y menciona ideas puntuales para trabajos futuros o mejoras del sistema \cite{reference1}.

% ------------------------------------------------
% Bibliografía
% ------------------------------------------------
\bibliographystyle{IEEEtran}
\bibliography{references} % Recuerda tener un archivo llamado 'references.bib' en tu proyecto

\end{document}
